% GENERATED FILE. Edit canonical lessons, then run npm run generate:books.
% canonical-source-hash: fnv1a64:baae69968ada63ff
% canonical-lessons: PA-C269-pachhanna, PA-C269-naksha, PA-C269-nishan, PA-C269-dhukkna, PA-C269-toa

\chapter{Finding the Way}
\label{ch:pa-a2-finding-the-way}
\hlchaptermodality{269}

\begin{chapteropening}
\textbf{By the end of this chapter:} \emph{I can find my way with a map and the signs.}

Complete the last lesson of chapter 269: I can find my way with a map and the signs.
\end{chapteropening}

\section[\texorpdfstring{\pa{ਪਛਾਣਨਾ} (pachhāṇnā)}{ਪਛਾਣਨਾ (pachhāṇnā)}]{\pa{ਪਛਾਣਨਾ} (pachhāṇnā) — to recognize}

\label{lesson:PA-C269-pachhanna}



\begin{quote}
Before the new one: say the Punjabi for a crowd, then the Punjabi for to stay put.
\end{quote}

\subsection*{You'll want to know: \pa{ਪਛਾਣਨਾ}}
\textbf{\pa{ਪਛਾਣਨਾ}} — \emph{pachhāṇnā} — \textquotedblleft{}to recognize\textquotedblright{}.

\begin{quote}
\pa{ਮੈਂ} \pa{ਇਹ} \pa{ਗਲੀ} \pa{ਪਛਾਣਦਾ} \pa{ਹਾਂ।}
\end{quote}

\begin{quote}
\emph{maiṁ eh galī pachhāṇdā hāṁ.}
\end{quote}

\begin{quote}
I recognize this lane.
\end{quote}

Past (I …): \pa{ਮੈਂ} \pa{ਰਸਤਾ} \pa{ਪਛਾਣਿਆ}, \emph{maiṁ rastā pachhāṇiā}.

\begin{cousinweb}
\textbf{\pa{ਪਛਾਣਨਾ}} (\emph{pachhāṇnā}) is inherited from Sanskrit \emph{pratyabhijānāti}, \textquotedblleft{}recognizes\textquotedblright{}; CDIAL lists Punjabi \emph{pachāṇnā}, \textquotedblleft{}recognizes\textquotedblright{}, under it (CDIAL 8637).
\end{cousinweb}

\begin{grammarlens}[title={What you've built}]
One more word for this chapter.
\end{grammarlens}

\subsection*{Your turn}
\begin{itemize}
\raggedright
  \item \emph{Say it:} \emph{pachhāṇnā}
  \item \emph{Say it:} \emph{pachhāṇnā}, once more
  \item \emph{Say it:} \emph{maiṁ eh galī pachhāṇdā hāṁ.} — I recognize this lane.
  \item \emph{Recall:} say the Punjabi for a crowd, then the Punjabi for to stay put, then say \emph{pachhāṇnā} again
\end{itemize}

\begin{tcolorbox}[breakable,colback=teal!4,colframe=teal!35!black,title={Before you move on}]
What does \pa{ਪਛਾਣਨਾ} mean? (\textquotedblleft{}To recognize\textquotedblright{}.) Say it once more.
\end{tcolorbox}
\section[\texorpdfstring{\pa{ਨਕਸ਼ਾ} (nakshā)}{ਨਕਸ਼ਾ (nakshā)}]{\pa{ਨਕਸ਼ਾ} (nakshā) — a map}

\label{lesson:PA-C269-naksha}



\begin{quote}
Before the new one: say the Punjabi for to stay put, then the Punjabi for to recognize.
\end{quote}

\subsection*{You'll want to know: \pa{ਨਕਸ਼ਾ}}
\textbf{\pa{ਨਕਸ਼ਾ}} — \emph{nakshā} — \textquotedblleft{}a map\textquotedblright{}.

\begin{quote}
\pa{ਮੇਰੇ} \pa{ਫ਼ੋਨ} \pa{ਵਿੱਚ} \pa{ਸ਼ਹਿਰ} \pa{ਦਾ} \pa{ਨਕਸ਼ਾ} \pa{ਹੈ।}
\end{quote}

\begin{quote}
\emph{mere fon vich shahir dā nakshā hai.}
\end{quote}

\begin{quote}
There is a map of the town on my phone.
\end{quote}

A masculine noun.

\begin{cousinweb}
\textbf{\pa{ਨਕਸ਼ਾ}} (\emph{nakshā}) is a loan, usually traced to Persian \emph{naqsha}, \textquotedblleft{}a map, a plan\textquotedblright{}, itself from Arabic.
\end{cousinweb}

\begin{grammarlens}[title={What you've built}]
One more word for this chapter.
\end{grammarlens}

\subsection*{Your turn}
\begin{itemize}
\raggedright
  \item \emph{Say it:} \emph{nakshā}
  \item \emph{Say it:} \emph{nakshā}, once more
  \item \emph{Say it:} \emph{mere fon vich shahir dā nakshā hai.} — There is a map of the town on my phone.
  \item \emph{Recall:} say the Punjabi for to stay put, then the Punjabi for to recognize, then say \emph{nakshā} again
  \item \emph{Write it:} copy \textbf{\pa{ਨਕਸ਼ਾ}} once, checking the whole word
\end{itemize}

\begin{tcolorbox}[breakable,colback=teal!4,colframe=teal!35!black,title={Before you move on}]
What does \pa{ਨਕਸ਼ਾ} mean? (\textquotedblleft{}A map\textquotedblright{}.) Say it once more.
\end{tcolorbox}
\section[\texorpdfstring{\pa{ਨਿਸ਼ਾਨ} (nishān)}{ਨਿਸ਼ਾਨ (nishān)}]{\pa{ਨਿਸ਼ਾਨ} (nishān) — a mark; a sign}

\label{lesson:PA-C269-nishan}



\begin{quote}
Before the new one: say the Punjabi for to recognize, then the Punjabi for a map.
\end{quote}

\subsection*{You'll want to know: \pa{ਨਿਸ਼ਾਨ}}
\textbf{\pa{ਨਿਸ਼ਾਨ}} — \emph{nishān} — \textquotedblleft{}a mark; a sign\textquotedblright{}.

\begin{quote}
\pa{ਸੜਕ} \pa{ਉੱਤੇ} \pa{ਲਾਲ} \pa{ਨਿਸ਼ਾਨ} \pa{ਹੈ।}
\end{quote}

\begin{quote}
\emph{saṛak utte lāl nishān hai.}
\end{quote}

\begin{quote}
There is a red mark on the road.
\end{quote}

A masculine noun.

\begin{cousinweb}
\textbf{\pa{ਨਿਸ਼ਾਨ}} (\emph{nishān}) is a loan, usually traced to Persian \emph{nishān}, \textquotedblleft{}a mark, a sign\textquotedblright{}.
\end{cousinweb}

\begin{grammarlens}[title={What you've built}]
One more word for this chapter.
\end{grammarlens}

\subsection*{Your turn}
\begin{itemize}
\raggedright
  \item \emph{Say it:} \emph{nishān}
  \item \emph{Say it:} \emph{nishān}, once more
  \item \emph{Say it:} \emph{saṛak utte lāl nishān hai.} — There is a red mark on the road.
  \item \emph{Recall:} say the Punjabi for to recognize, then the Punjabi for a map, then say \emph{nishān} again
  \item \emph{Write it:} copy \textbf{\pa{ਨਿਸ਼ਾਨ}} once, checking the whole word
\end{itemize}

\begin{tcolorbox}[breakable,colback=teal!4,colframe=teal!35!black,title={Before you move on}]
What does \pa{ਨਿਸ਼ਾਨ} mean? (\textquotedblleft{}A mark; a sign\textquotedblright{}.) Say it once more.
\end{tcolorbox}
\section[\texorpdfstring{\pa{ਢੁੱਕਣਾ} (ḍhukkṇā)}{ਢੁੱਕਣਾ (ḍhukkṇā)}]{\pa{ਢੁੱਕਣਾ} (ḍhukkṇā) — to draw near; to approach}

\label{lesson:PA-C269-dhukkna}



\begin{quote}
Before the new one: say the Punjabi for a map, then the Punjabi for a mark.
\end{quote}

\subsection*{You'll want to know: \pa{ਢੁੱਕਣਾ}}
\textbf{\pa{ਢੁੱਕਣਾ}} — \emph{ḍhukkṇā} — \textquotedblleft{}to draw near; to approach\textquotedblright{}.

\begin{quote}
\pa{ਬੱਸ} \pa{ਹੌਲੀ} \pa{ਹੌਲੀ} \pa{ਅੱਡੇ} \pa{ਦੇ} \pa{ਨੇੜੇ} \pa{ਢੁੱਕਦੀ} \pa{ਹੈ।}
\end{quote}

\begin{quote}
\emph{bass haulī haulī aḍḍe de neṛe ḍhukkdī hai.}
\end{quote}

\begin{quote}
The bus slowly draws near the stand.
\end{quote}

Past (it …): \pa{ਉਹ} \pa{ਬੱਸ} \pa{ਅੱਡੇ} \pa{ਦੇ} \pa{ਨੇੜੇ} \pa{ਢੁੱਕੀ}, \emph{oh bass aḍḍe de neṛe ḍhukkī}.

\begin{cousinweb}
\textbf{\pa{ਢੁੱਕਣਾ}} (\emph{ḍhukkṇā}) is inherited from a reconstructed \emph{*ḍhukyati}, \textquotedblleft{}approaches\textquotedblright{}; CDIAL lists Punjabi \emph{ḍhukkṇā}, \textquotedblleft{}to approach\textquotedblright{}, under it (CDIAL 5592).
\end{cousinweb}

\begin{grammarlens}[title={What you've built}]
One more word for this chapter.
\end{grammarlens}

\subsection*{Your turn}
\begin{itemize}
\raggedright
  \item \emph{Say it:} \emph{ḍhukkṇā}
  \item \emph{Say it:} \emph{ḍhukkṇā}, once more
  \item \emph{Say it:} \emph{bass haulī haulī aḍḍe de neṛe ḍhukkdī hai.} — The bus slowly draws near the stand.
  \item \emph{Recall:} say the Punjabi for a map, then the Punjabi for a mark, then say \emph{ḍhukkṇā} again
  \item \emph{Write it:} copy \textbf{\pa{ਢੁੱਕਣਾ}} once, checking the whole word
\end{itemize}

\begin{tcolorbox}[breakable,colback=teal!4,colframe=teal!35!black,title={Before you move on}]
What does \pa{ਢੁੱਕਣਾ} mean? (\textquotedblleft{}To draw near; to approach\textquotedblright{}.) Say it once more.
\end{tcolorbox}
\section[\texorpdfstring{\pa{ਟੋਆ} (ṭoā)}{ਟੋਆ (ṭoā)}]{\pa{ਟੋਆ} (ṭoā) — a pothole; a hole in the ground}

\label{lesson:PA-C269-toa}



\begin{quote}
Before the new one: say the Punjabi for a mark, then the Punjabi for to draw near.
\end{quote}

\subsection*{You'll want to know: \pa{ਟੋਆ}}
\textbf{\pa{ਟੋਆ}} — \emph{ṭoā} — \textquotedblleft{}a pothole; a hole in the ground\textquotedblright{}.

\begin{quote}
\pa{ਸੜਕ} \pa{ਵਿੱਚ} \pa{ਇੱਕ} \pa{ਡੂੰਘਾ} \pa{ਟੋਆ} \pa{ਹੈ।}
\end{quote}

\begin{quote}
\emph{saṛak vich ikk ḍū̃ghā ṭoā hai.}
\end{quote}

\begin{quote}
There is a deep pothole in the road.
\end{quote}

A masculine noun.

\begin{cousinweb}
\textbf{\pa{ਟੋਆ}} (\emph{ṭoā}) is inherited from a reconstructed \emph{*trōpa}, \textquotedblleft{}piercing, hole\textquotedblright{}; CDIAL lists Punjabi \emph{ṭoā}, \textquotedblleft{}hole, pit\textquotedblright{}, under it (CDIAL 6082).
\end{cousinweb}

\begin{grammarlens}[title={What you've built}]
One more word for this chapter.
\end{grammarlens}

\subsection*{Your turn}
\begin{itemize}
\raggedright
  \item \emph{Say it:} \emph{ṭoā}
  \item \emph{Say it:} \emph{ṭoā}, once more
  \item \emph{Say it:} \emph{saṛak vich ikk ḍū̃ghā ṭoā hai.} — There is a deep pothole in the road.
  \item \emph{Recall:} say the Punjabi for a mark, then the Punjabi for to draw near, then say \emph{ṭoā} again
  \item \emph{Write it:} copy \textbf{\pa{ਟੋਆ}} once, checking the whole word
\end{itemize}

\begin{tcolorbox}[breakable,colback=teal!4,colframe=teal!35!black,title={Before you move on}]
What does \pa{ਟੋਆ} mean? (\textquotedblleft{}A pothole; a hole in the ground\textquotedblright{}.) Say it once more.
\end{tcolorbox}
